\documentclass[12pt]{article}
\usepackage[utf8]{inputenc}
\usepackage{amsmath}
\usepackage{amssymb}
\usepackage{bm}

\title{SEU鲁棒训练目标函数}
\author{}
\date{}

\begin{document}

\maketitle

\section{统一SEU鲁棒目标函数}

以SEU容错为出发点的训练目标函数：

\begin{equation}
\min_{\theta}\;\mathbb{E}_{(x,y)\sim\mathcal{D}}
\Big[
\underbrace{\ell\!\left(f(x;\,W_\theta),y\right)}_{\text{干净任务损失}}
\;+\;
\lambda_{\text{rob}}\underbrace{\mathbb{E}_{\xi\sim\mathcal{D}_{\text{SEU}}}\Big[\ell\!\left(f(x;\,W_\theta+\Delta W(\xi)),y\right)\Big]}_{\text{SEU期望风险（Monte Carlo/多次采样）}}
\;+\;
\lambda_{\text{qe}}\underbrace{\sum_{i}\!\left[\big(\mathrm{clip}(w_i) - q(w_i)\big)^2\right]}_{\text{量化误差最小化（QE Loss）}}
\;+\;
\lambda_{\text{sens}}\underbrace{\sum_{i} s_i\cdot \mathbb{E}_{\xi}\!\left[\big(\Delta q_i(\xi)\big)^2\right]}_{\text{敏感度加权的编码扰动惩罚}}
\;+\;
\lambda_{\text{wd}}\underbrace{\|W_\theta\|_2^2}_{\text{幅度约束/分布收缩}}
\Big]
\label{eq:seu_robust}
\end{equation}

\section{符号说明}

\begin{itemize}
    \item $f(x;W_\theta)$：神经网络前向传播函数，$W_\theta$为可训练权重参数
    \item $p_\theta(x)=\mathrm{softmax}(f(x;W_\theta))$：网络输出的概率分布
    \item $\ell(\cdot,\cdot)$：损失函数（如交叉熵）
    \item $\mathcal{D}$：训练数据分布
    \item $\xi \sim \mathcal{D}_{\text{SEU}}$：按BER（Bit Error Rate）采样的bit-flip随机变量，可包含层/bit的非均匀概率分布
    \item $\Delta W(\xi)$：由bit-flip诱导的权重扰动，通常满足$\Delta W = s\cdot \Delta q$，其中$\Delta q$是整数量化码的变化
    \item $s_i$：第$i$个权重量化码的"重要性/敏感度"，可取$s_i=\mathbb{E}[(\partial \ell/\partial w_i)^2]$（梯度平方的期望，对应之前讨论的梯度平方敏感度）
    \item $\Delta q_i(\xi)$：第$i$个量化码在故障$\xi$下的变化量
    \item $\mathrm{clip}(w_i)$：第$i$个权重被裁剪到量化范围内的值
    \item $q(w_i)$：第$i$个权重的量化值
    \item $\lambda_{\text{rob}}$：SEU期望风险的权重系数
    \item $\lambda_{\text{qe}}$：量化误差最小化的权重系数（对应代码中的\texttt{QE\_loss\_weight}）
    \item $\lambda_{\text{sens}}$：敏感度加权编码扰动惩罚的权重系数
    \item $\lambda_{\text{wd}}$：权重衰减（L2正则化）的权重系数
\end{itemize}

\section{与现有实现机制的对应关系}

\begin{itemize}
    \item \textbf{$\mathbb{E}_{\xi\sim\mathcal{D}_{\text{SEU}}}[\cdot]$的Monte Carlo估计}：对应代码中的\texttt{nr\_random\_sample}机制，通过多次随机采样bit-flip故障来估计期望风险
    
    \item \textbf{$\sum_i s_i\cdot \mathbb{E}_{\xi}[(\Delta q_i)^2]$项}：将QE Loss和distribution loss重新解释为"SEU误差传播上界"的叙事角度。具体而言：
    \begin{itemize}
        \item 敏感度$s_i$高的权重，其量化码翻转的影响更大
        \item 通过惩罚$\mathbb{E}_{\xi}[(\Delta q_i)^2]$，鼓励量化码在故障下保持稳定
        \item 这等价于最小化敏感度加权的量化误差（QE Loss的SEU视角）
        \item 同时通过权重衰减$\|W_\theta\|_2^2$，收缩权重分布，减少故障传播幅度（distribution loss的SEU视角）
    \end{itemize}
    
    \item \textbf{量化误差最小化项（QE Loss）}：最小化权重与其量化值之间的平方误差，即$\sum_i[\mathrm{clip}(w_i) - q(w_i)]^2$。从SEU视角看，这鼓励权重更接近量化网格点，从而减少bit-flip导致的量化码变化幅度，降低故障影响。对应代码中的\texttt{auxiliary\_quantized\_loss}函数计算的\texttt{QE\_loss}
\end{itemize}

\section{实现建议}

\begin{enumerate}
    \item \textbf{替换位宽采样为SEU采样}：将\texttt{nr\_random\_sample}中的位宽配置采样改为按BER采样bit-flip故障$\xi$，每次前向传播时应用$\Delta W(\xi)$
    
    \item \textbf{敏感度计算}：在训练过程中收集梯度统计信息，计算$s_i=\mathbb{E}[(\partial \ell/\partial w_i)^2]$，用于加权编码扰动惩罚项
    
    \item \textbf{Monte Carlo估计}：通过多次采样$\xi$并计算平均损失来估计$\mathbb{E}_{\xi\sim\mathcal{D}_{\text{SEU}}}[\ell(\cdot)]$
    
    \item \textbf{超参数调优}：根据实验效果调整$\lambda_{\text{rob}}$、$\lambda_{\text{qe}}$、$\lambda_{\text{sens}}$、$\lambda_{\text{wd}}$的权重
\end{enumerate}

\end{document}

